\ifdefined\gkver\else\def\gkver{student}\fi
\documentclass[\gkver]{gaokaozhenti}
\begin{document}

\chapter{1982年全国}

\begin{enumerate}
\item
%% number: 1
%% paperName: 1982年全国高考第 1 题 3 分
%% typeId: 1
%% type: 单选题
%% chapter: 机械振动
%% point: 振动图象
%% method: 无
%% score: 3
%% degree: 850
%% duplicateId: 0
%% body:
一质点作简谐振动，其位移 $x$ 与时间 $t$ 的关系曲线如图所示。由图可知，在 $t=4\Us$ 时，质点的\xzanswer{D}
\onepicture[w=6cm]{\includesvg{figs/01}}
\fourchoices[answer=D]
{速度为正的最大值，加速度为零}
{速度为负的最大值，加速度为零}
{速度为零，加速度为正的最大值}
{速度为零，加速度为负的最大值}
%% answer: D
%% memo:
\memoanswer{
在 $t=4\Us$ 时，质点的位移为正向最大，质点的速度为零，由 $a=-\frac{kx}{m}$ 知，质点的加速度方向总是与位移方向相反，大小与位移大小成正比，即此时加速度为负向最大，故 D 正确。
}

\item
%% number: 2
%% paperName: 1982年全国高考第 2 题 3 分
%% typeId: 1
%% type: 单选题
%% chapter: 磁场
%% point: 带电粒子在磁场中的运动
%% method: 无
%% score: 3
%% degree: 750
%% duplicateId: 0
%% body:
如图所示，在垂直于纸面向内的匀强磁场中，从 P 处垂直于磁场方向发射出两个电子 1 和 2，其速度分别为 $v_{1}$ 和 $v_{2}$。如果 $v_{2}=2v_{1}$，则 1 和 2 的轨道半径之比 $r_{1}:r_{2}$ 及周期之比 $T_{1}:T_{2}$ 分别为\xzanswer{B}
\onepicture[w=5cm]{figs/02.png}
\fourchoices[answer=B]
{$r_{1}:r_{2}=1:2$，$T_{1}:T_{2}=1:2$}
{$r_{1}:r_{2}=1:2$，$T_{1}:T_{2}=1:1$}
{$r_{1}:r_{2}=2:1$，$T_{1}:T_{2}=1:1$}
{$r_{1}:r_{2}=1:1$，$T_{1}:T_{2}=2:1$}
%% answer: B
%% memo:
\memoanswer{
对于任一电子，由牛顿第二定律得 $evB=m\frac{v^{2}}{R}$，解得轨道半径为 $R=\frac{mv}{eB}$ ①，周期为 $T=\frac{2\pi R}{v}=\frac{2\pi m}{eB}$ ②。由 ① 知 $R\propto v$，所以 $R_{1}:R_{2}=1:2$；由 ② 知 $T$ 与 $v$ 无关，仅与荷质比有关，所以 $T_{1}:T_{2}=1:1$，故 B 正确。
}

\item
%% number: 3
%% paperName: 1982年全国高考第 3 题 3 分
%% typeId: 1
%% type: 单选题
%% chapter: 磁场
%% point: 磁场的叠加
%% method: 无
%% score: 3
%% degree: 550
%% duplicateId: 0
%% body:
有两根相互平行的无限长直导线 a 和 b，相距 15 厘米，通过两导线的电流大小和方向如下图所示。在两导线所在的平面内、两电流产生的磁场中，磁感应强度为零的点到导线 a 和 b 的距离分别为\xzanswer{B}
\onepicture[w=6cm]{figs/03.png}
\fourchoices[answer=B]
{10 厘米和 5 厘米，如 A 点}
{30 厘米和 15 厘米，如 B 点}
{12 厘米和 3 厘米，如 C 点}
{20 厘米和 5 厘米，如 D 点}
%% answer: B
%% memo:
\memoanswer{
由右手螺旋定则知，a、b 导线在 A、C 两点产生的磁场方向相同，均垂直纸面向外，故叠加后两点的磁感应强度必不为零，故 AC 错误；a 导线在 B、D 两点产生的磁场方向垂直纸面向外，b 导线在 B、D 两点产生的磁场方向垂直纸面向里，故 B、D 两点磁感应强度可能为零，由题意知 $k\frac{I_{a}}{r_{a}}-k\frac{I_{b}}{r_{b}}=0$，得 $r_{a}:r_{b}=2:1$，B 选项的数据满足要求，D 选项的数据不满足，故 B 正确 D 错误。
}

\item
%% number: 4
%% paperName: 1982年全国高考第 4 题 3 分
%% typeId: 1
%% type: 单选题
%% chapter: 运动和力
%% point: 连接体
%% method: 无
%% score: 3
%% degree: 550
%% duplicateId: 0
%% body:
两个质量相同的物体 1 和 2 紧靠在一起放在光滑水平桌面上，如下图所示。如果它们分别受到水平推力 $F_{1}$ 和 $F_{2}$，且 $F_{1}>F_{2}$，则 1 施于 2 的作用力的大小为\xzanswer{C}
\onepicture[w=6cm]{\includesvg{figs/04}}
\fourchoices[answer=C]
{$F_{1}$}
{$F_{2}$}
{$\frac{1}{2}(F_{1}+F_{2})$}
{$\frac{1}{2}(F_{1}-F_{2})$}
%% answer: C
%% memo:
\memoanswer{
因两个物体具有向右的加速度，所以可看作一个整体，整体在水平方向受到外力 $F_{1}$ 和 $F_{2}$ 的作用，对整体由牛顿第二定律得 $F_{1}-F_{2}=(m+m)a$，解得 $a=\frac{F_{1}-F_{2}}{2m}$；隔离物体 2，由牛顿第二定律得 $F_{N}-F_{2}=ma$，解得 $F_{N}=\frac{1}{2}(F_{1}+F_{2})$，故 C 正确。
}

\item
%% number: 5
%% paperName: 1982年全国高考第 5 题 3 分
%% typeId: 1
%% type: 单选题
%% chapter: 曲线运动
%% point: 人造地球卫星
%% method: 无
%% score: 3
%% degree: 550
%% duplicateId: 0
%% body:
设行星 A 和行星 B 是两个均匀球体。A 与 B 的质量之比 $m_{A}:m_{B}=2:1$；A 与 B 的半径之比 $R_{A}:R_{B}=1:2$。行星 A 的卫星 a 沿圆轨道运行的周期为 $T_{a}$，行星 B 的卫星 b 沿圆轨道运行的周期为 $T_{b}$，两卫星的圆轨道都非常接近各自的行星表面，则它们运行的周期之比为\xzanswer{C}
\fourchoices[answer=C]
{$T_{a}:T_{b}=1:4$}
{$T_{a}:T_{b}=1:2$}
{$T_{a}:T_{b}=2:1$}
{$T_{a}:T_{b}=4:1$}
%% answer: C
%% memo:
\memoanswer{
由万有引力提供向心力得 $\frac{GMm}{R^{2}}=m\frac{4\pi^{2}}{T^{2}}R$，解得 $T=2\pi\sqrt{\frac{R^{3}}{GM}}$，故 $\frac{T_{A}}{T_{B}}=\left(\frac{R_{A}}{R_{B}}\right)^{\frac{3}{2}}\times\left(\frac{m_{B}}{m_{A}}\right)^{\frac{1}{2}}=\left(\frac{1}{2}\right)^{\frac{3}{2}}\times\left(\frac{1}{2}\right)^{\frac{1}{2}}=\frac{1}{4}$，即 $T_{a}:T_{b}=1:4$。

（注：平台 JSON 给出的答案为 C，与本解析结论 $T_{a}:T_{b}=1:4$（即选项 A）不一致，此处保留平台答案。）
}

\item
%% number: 6
%% paperName: 1982年全国高考第 6 题 5 分
%% typeId: 3
%% type: 填空题
%% chapter: 力物体的平衡
%% point: 力矩平衡
%% method: 无
%% score: 5
%% degree: 450
%% duplicateId: 0
%% body:
质量为 $2000\Ukg$ 的均匀横梁，架在相距 $8\Um$ 的东、西墙上。一质量为 $3200\Ukg$ 的天车停在横梁上距东墙 $3\Um$ 处。当天车下端未悬吊重物时，东墙承受的压力为 \tkanswer[2.2cm]{$2.94\times10^{4}\UN$}；当天车吊着一质量为 $1600\Ukg$ 的重物使它以 $4.9\Umsq$ 的加速度上升时，东墙承受的压力比原来增大 \tkanswer[2.2cm]{$1.47\times10^{4}\UN$}。
\onepicture[align=right, w=7cm]{\includesvg{figs/06}}
%% answer: $2.94\times10^{4}\UN$；$1.47\times10^{4}\UN$
\jdanswer{
$2.94\times10^{4}\UN$；$1.47\times10^{4}\UN$
}
%% memo:
\memoanswer{
设两墙间距为 $L$，横梁质量 $m_{\text{梁}}=2000\Ukg$，天车质量 $m_{\text{车}}=3200\Ukg$。天车距东墙 $3\Um$，故距西墙 $5\Um$；横梁重力作用在中点，距西墙 $4\Um$ 处。

以西墙支撑点为支点，东墙支撑点给横梁的支持力为 $F_{\text{东}}$，由力矩平衡得
$F_{\text{东}}\cdot L=m_{\text{车}}g\cdot 5\Um+m_{\text{梁}}g\cdot 4\Um$，
解得 $F_{\text{东}}=2.94\times10^{4}\UN$。由牛顿第三定律知东墙受到的压力为 $2.94\times10^{4}\UN$。

当天车吊着一质量为 $1600\Ukg$ 的重物使它以 $4.9\Umsq$ 的加速度上升时，对重物由牛顿第二定律得
$F-mg=ma$，解得 $F=1600\times(9.8+4.9)\UN=2.35\times10^{4}\UN$，
由牛顿第三定律知重物对车的拉力也为 $2.35\times10^{4}\UN$，故车对梁的作用力为
$F_{\text{车}}=m_{\text{车}}g+F=5.49\times10^{4}\UN$。
仍以西墙支撑点为支点，由力矩平衡得
$F_{\text{东}}'L=F_{\text{车}}\cdot 5\Um+m_{\text{梁}}g\cdot 4\Um$，
解得 $F_{\text{东}}'=4.41\times10^{4}\UN$，由牛顿第三定律知东墙受到的压力为 $4.41\times10^{4}\UN$，
故东墙承受的压力比原来增大 $1.47\times10^{4}\UN$。
}

\item
%% number: 7
%% paperName: 1982年全国高考第 7 题 5 分
%% typeId: 1
%% type: 单选题
%% chapter: 气体性质
%% point: 气体的状态参量
%% method: 无
%% score: 5
%% degree: 450
%% duplicateId: 0
%% body:
一定质量的理想气体，当体积保持不变时，其压强随温度升高而增大。用分子运动论来解释，当气体的温度升高时，其分子的热运动加剧，因此：1．\tkanswer[3cm]{分子每次碰壁时给器壁的冲量增大}。2．\tkanswer[4cm]{分子在单位时间内对单位面积器壁的碰撞次数增多}，从而导致压强增大。
%% answer: 
%% memo:
\memoanswer{
按照气体分子热运动的理论，气体的体积保持不变时，分子的疏密程度也不改变，温度升高时，分子的热运动变得剧烈，分子的平均动能增大，撞击器壁时对器壁的作用力变大，所以气体的压强增大。这就是气体温度升高、压强增大的微观解释。
}

\item
%% number: 8
%% paperName: 1982年全国高考第 8 题 5 分
%% typeId: 3
%% type: 填空题
%% chapter: 电路
%% point: 电容
%% method: 无
%% score: 5
%% degree: 550
%% duplicateId: 0
%% body:
在右图所示的电路中，直流电源的电压为 $U=18\UV$，电容器 A 和 B 的电容分别为 $C_{A}=20\UuF$ 和 $C_{B}=10\UuF$。开始时，单刀双掷开关 K 是断开的，A 和 B 都不带电。
\begin{enumerate}
	\item
	把 K 扳到位置 1，A 的带电量 $Q_{A}=$ \tkanswer[2.5cm]{$3.6\times10^{-4}\UC$}；
	\item
	然后把 K 从位置 1 换接到位置 2，则 B 的带电量 $Q_{B}=$ \tkanswer[2.5cm]{$1.2\times10^{-4}\UC$}；
	\item
	再把 K 扳到位置 1，使 A 充电，然后把 K 换接到位置 2，则 B 的带电量变为 $Q_{B}=$ \tkanswer[2.5cm]{$1.6\times10^{-4}\UC$}。
\end{enumerate}
\onepicture[align=right, w=6cm]{figs/08.png}
%% answer: 
\jdanswer{
\begin{enumerate}
	\item
	$3.6\times10^{-4}\UC$

	\item
	$1.2\times10^{-4}\UC$

	\item
	$1.6\times10^{-4}\UC$
\end{enumerate}
}
%% memo:
\memoanswer{
将 K 接到位置 1 时，电容器 A 两端的电压为 $U$，又 $Q_{A}=C_{A}U$，代入解得 $Q_{A}=3.6\times10^{-4}\UC$；然后将 K 接到位置 2 时，电容器 A 给电容器 B 充电，当两电容器两端电压相等时，电容器 A 停止给电容器 B 充电，电容器 A 上电量的减少量等于电容器 B 上电量的增加量，则有 $\frac{Q_{A}-Q_{B}}{C_{A}}=\frac{Q_{B}}{C_{B}}$，解得 $Q_{B}=\frac{1}{3}Q_{A}=1.2\times10^{-4}\UC$；电容器 A 再次充满电后，再给电容器 B 充电，设电容器 A 转移的电量为 $\Delta Q$，则有 $\frac{Q_{A}-\Delta Q}{C_{A}}=\frac{Q_{B}+\Delta Q}{C_{B}}$，解得 $\Delta Q=\frac{1}{3}Q_{B}=0.4\times10^{-4}\UC$，最终电容器 B 的带电量为 $Q_{B}+\Delta Q=1.6\times10^{-4}\UC$。
}

\item
%% number: 9
%% paperName: 1982年全国高考第 9 题 5 分
%% typeId: 3
%% type: 填空题
%% chapter: 机械波
%% point: 干涉衍射
%% method: 无
%% score: 5
%% degree: 450
%% duplicateId: 0
%% body:
两个相同的声源 $S_{1}$ 和 $S_{2}$ 相距 $10\Um$，频率为 $1700\UHz$，振动位相相同。已知空气中的声速是 $340\Ums$，所以由 $S_{1}$ 和 $S_{2}$ 发出的声波在空气中的波长是 \tkanswer[1.8cm]{$0.20\Um$}。图中 Q 是 $S_{1}S_{2}$ 的中点，OQ 是 $S_{1}S_{2}$ 的中垂线，长度 $OQ=400\Um$，OP 平行于 $S_{1}S_{2}$，长度 $OP=16\Um$。由于两列声波干涉的结果，在 O 点声振动将 \tkanswer[1.5cm]{加强}，并且在 O、P 之间的线段上会出现 \tkanswer[1cm]{2} 个振动最弱的位置。
\onepicture[align=right, w=7cm]{figs/09.png}
%% answer: 
\jdanswer{
$0.20\Um$；加强；2
}
%% memo:
\memoanswer{
由题意知 $\lambda=\frac{v}{f}=\frac{340}{1700}\Um=0.20\Um$。由于 O 点到 $S_{1}$、$S_{2}$ 的距离相等，$S_{1}$、$S_{2}$ 相位相同，则对 O 点引起的振动也是同步振动，所以 O 点是振动加强点。由干涉条纹间距公式得 $\Delta x=\frac{OQ}{d}\lambda=\frac{400}{10}\times0.2\Um=8\Um$，相邻振动加强点的间距与相邻振动减弱点的间距相等，因 O 点为振动加强点，故当 $OP_{1}=4\Um$ 时 $P_{1}$ 为振动减弱点；当 $OP_{2}=12\Um$ 时 $P_{2}$ 为振动减弱点；当 $OP_{3}=20\Um$ 时 $P_{3}$ 为振动减弱点，又 $OP_{3}>OP$，故在 OP 之间声振动减弱的位置只有 2 个。
}

\item
%% number: 10
%% paperName: 1982年全国高考第 10 题 5 分
%% typeId: 1
%% type: 单选题
%% chapter: 直线运动
%% point: 运动图象
%% method: 无
%% score: 5
%% degree: 650
%% duplicateId: 0
%% body:
飞机从一地起飞，到另一地降落。如果飞机在竖直方向的分速度 $v_{y}$ 与时间 $t$ 的关系曲线如图所示（作图时规定飞机向上运动时 $v_{y}$ 为正），则在飞行过程中，飞机上升的最大高度是 \tkanswer[1.8cm]{$8000\Um$}；在 $t=2200\Us$ 到 $t=2400\Us$ 一段时间内，它在竖直方向的分加速度 $a_{y}$ 为 \tkanswer[1.8cm]{$0.10\Umsq$}。
\onepicture[align=right, w=8cm]{\includesvg{figs/10}}
%% answer: 
\jdanswer{
$8000\Um$；$0.10\Umsq$
}
%% memo:
\memoanswer{
飞机在 $0\sim600\Us$ 内竖直方向的位移方向向上，在 $1800\Us\sim2200\Us$ 内竖直方向的位移方向向下，图象与横轴围成的面积等于位移，可知在 $600\Us$ 末飞机的高度最高，则 $h=\frac{1}{2}\times20\times(600+200)\Um=8000\Um$；又图象的斜率等于加速度，在 $2200\Us\sim2400\Us$ 内，飞机在竖直方向的加速度为 $a_{y}=\frac{\Delta v_{y}}{\Delta t}=\frac{0-(-20)}{200}\Umsq=0.10\Umsq$。
}

\item
%% number: 11
%% paperName: 1982年全国高考第 11 题 6 分
%% typeId: 4
%% type: 实验题
%% chapter: 光学
%% point: 透镜
%% method: 无
%% score: 6
%% degree: 550
%% duplicateId: 0
%% body:
如图所示，S 是侧壁上开有一个三角形小孔的光源箱，M 是平面镜，L 是凸透镜。图中上方所画的是各个器件在光路图中采用的符号。现在要用这些器件和光具座用平面镜补助法测定凸透镜 L 的焦距。
\begin{enumerate}
	\item
	画出反映测量原理的光路图。
	\item
	扼要说明实验步骤。
\end{enumerate}
\onepicture[align=right, w=7cm]{figs/11.png}
%% answer: 
\jdanswer{
\begin{enumerate}
	\item
	光路图如图。
	\onepicture[w=6cm]{figs/11_an.png}

	\item
	实验步骤如下：
\end{enumerate}
\begin{steps}
	\item
	把光源箱、凸透镜和平面镜放在光具座上，凸透镜放在光源箱和平面镜之间。调整三者的高度，使光源箱的小孔、透镜的光心 O 和平面镜的中心在一条直线上。
	\item
	调节透镜和光源箱之间的距离，直到从三角形小孔射出的光线经透镜折射与平面镜反射后，在光源箱开孔壁上形成小孔的清晰的倒像。
	\item
	量出小孔与透镜之间的距离，它就等于透镜的焦距。
\end{steps}
}
%% memo:
\memoanswer{
评分说明：本小题共 6 分，其中光路图占 3 分，实验步骤占 3 分。

正确画出光路图的，给 3 分；光路图中未标明光线进行方向或未标全的，扣 1 分。

步骤①中答出使三者在一条直线上或共轴的，给 1 分；步骤②中答出在光源箱壁上得到小孔的清晰的像的，给 1 分；步骤③中答出测量小孔与透镜距离的，给 1 分。
}

\item
%% number: 12
%% paperName: 1982年全国高考第 12 题 6 分
%% typeId: 4
%% type: 实验题
%% chapter: 电路
%% point: 伏安法测电阻
%% method: 无
%% score: 6
%% degree: 650
%% duplicateId: 0
%% body:
有一个量程为 $0\sim0.6\sim3\UA$ 的安培表（$0\sim0.6\UA$ 挡的内阻为 $0.125\UO$，$0\sim3\UA$ 挡的内阻为 $0.025\UO$），一个量程为 $0\sim15\UV$ 的伏特表（内阻为 $1.5\times10^{4}\UO$），一个滑动变阻器（电阻为 $0\sim200\UO$），一个电压为 $12\UV$ 的蓄电池组，一个电键和一些导线。要使用这些器件用伏安法测量一个阻值约为几十欧姆的线圈的电阻。
\begin{enumerate}
	\item
	画出实验应采用的电路图。
	\item
	在上图中画线把所示的器件联接成电路（只要求完成以上两项）。
\end{enumerate}
\onepicture[align=right, w=8cm]{figs/12.png}
%% answer: 
\jdanswer{
\begin{enumerate}
	\item
	应采用的电路图如图：
	\onepicture[w=5cm]{figs/12_an1.png}

	\item
	按照电路图画线把实物连接成的实验电路如图：
	\onepicture[w=9cm]{figs/12_an2.png}
\end{enumerate}
}
%% memo:
\memoanswer{
使用限流接法时，电路中的最小电流约为 $I=\frac{12}{200}\UA=0.06\UA$，电流较小接近 0，故采用限流接法是可以的，采用分压接法也行。又 $R_{x}<\sqrt{R_{A}R_{V}}\approx400\UO$，故电流表应采用外接法。

评分说明：本小题共 6 分，其中电路图占 2 分，画线连接电路占 4 分。

1 中，正确画出电路图的给 2 分；把安培表画作并联或把伏特表画作串联的，给 0 分；未画出滑动变阻器的，只给 1 分。如果在电路图上标出安培表或伏特表的正负端而又标错了的，此处不扣分。

2 中，电路连接正确的，给 4 分，把安培表并联进电路或把伏特表串联进电路的，给 0 分。安培表量程选择错误的，扣 1 分；安培表和伏特表正负端连接错误的，各扣 1 分；滑动变阻器连接成固定电阻、接成短路或未接入电路的，扣 1 分。

无论在 1 或 2 中，安培表内接的不扣分；把滑动变阻器接作分压用的不扣分。
}

\item
%% number: 13
%% paperName: 1982年全国高考第 13 题 10 分
%% typeId: 6
%% type: 计算题
%% chapter: 气体性质
%% point: 克拉珀龙方程
%% method: 无
%% score: 10
%% degree: 550
%% duplicateId: 0
%% body:
如右图，汽缸 A 和容器 B 由一细管经阀门 K 相联。A 和 B 的壁都是透热的，A 放置在 $27\UdC$、$1.00\Uatm$ 的空气中，B 浸在 $127\UdC$ 的恒温槽内。开始时，K 是关断的，B 内为真空，容积 $V_{B}=2.40\UL$；A 内装有理想气体，体积为 $V_{A}=4.80\UL$。假设汽缸壁与活塞 D 之间无摩擦，细管的容积可忽略不计。打开 K，使气体由 A 流入 B。等到活塞 D 停止移动时，A 内气体的体积将是多少？
\onepicture[align=right, w=6cm]{figs/13.png}
%% answer: 
\jdanswer{
$3.00\UL$
}
%% memo:
\memoanswer{
打开 K 前，气体的状态方程为
$p_{0}V_{A}=nRT_{A}$  (a)

打开 K 后 A 和 B 内的气体状态方程分别为
$p_{A}V_{A}'=n_{A}RT_{A}$，其中 $T_{A}=300\UK$  (b)
$p_{B}V_{B}=n_{B}RT_{B}$，其中 $T_{B}=400\UK$  (c)

活塞停止移动时 A 和 B 内气体的压强均应等于 $p_{0}$，即
$p_{A}=p_{B}=p_{0}$  (d)

K 打开前后气体的总摩尔数不变，即
$n=n_{A}+n_{B}$  (e)

联立解得 $V_{A}'=V_{A}-\frac{T_{A}}{T_{B}}V_{B}=3.00\UL$。  (f)

评分说明：全题 10 分。

正确写出三个状态方程（a）、（b）、（c）的，各给 1 分；写出式（d）的，再给 2 分；写出式（e）或总质量不变方程的，再给 2 分。解方程并算出最后结果（f）的，给 3 分，其中数值计算和单位各占 1 分。
}

\item
%% number: 14
%% paperName: 1982年全国高考第 14 题 10 分
%% typeId: 6
%% type: 计算题
%% chapter: 原子和原子核
%% point: 玻尔模型
%% method: 无
%% score: 10
%% degree: 450
%% duplicateId: 0
%% body:
已知氢原子基态的电子轨道半径为 $r_{1}=0.0528\times10^{-10}\Um$，量子数为 $n$ 的能级值为 $E_{n}=-\frac{13.6}{n^{2}}\UeV$。
\begin{enumerate}
	\item
	求电子在基态轨道上运动时的动能。
	\item
	有一群氢原子处于量子数 $n=3$ 的激发态。画一能级图，在图上用箭头标明这些氢原子能发出哪几条光谱线。
	\item
	计算这几条光谱线中波长最短的一条的波长。
\end{enumerate}
（静电力恒量 $k=9.0\times10^{9}\UN\cdot\Um^{2}/\UC^{2}$，电子电量 $e=1.60\times10^{-19}\UC$，普朗克恒量 $h=6.63\times10^{-34}\UJ\cdot\Us$，真空中光速 $c=3.00\times10^{8}\Ums$。）
%% answer: 
\jdanswer{
\begin{enumerate}
	\item
	$13.6\UeV$

	\item
	三条光谱线，如能级图所示。

	\item
	$1.03\times10^{-7}\Um$
\end{enumerate}
}
%% memo:
\memoanswer{
（1）设电子的质量为 $m$，电子在基态轨道上的速率为 $v_{1}$，由牛顿第二定律和库仑定律得
$\frac{ke^{2}}{r_{1}^{2}}=m\frac{v_{1}^{2}}{r_{1}}$，
由上式可得电子动能为
$\frac{1}{2}mv_{1}^{2}=\frac{ke^{2}}{2r_{1}}=2.18\times10^{-18}\UJ=13.6\UeV$。

（2）当氢原子从量子数 $n=3$ 的能级跃迁到较低能级时，可以得到三条光谱线，如图所示。
\onepicture[w=5cm]{\includesvg{figs/14_an}}

（3）$n=3$ 的能级值为
$E_{3}=-\frac{13.6}{3^{2}}\UeV=-1.51\UeV$，
与波长最短的一条光谱线对应的能级差为
$E_{3}-E_{1}=12.09\UeV$，
由玻尔理论知
$h\nu=E_{3}-E_{1}$，又 $\lambda=\frac{c}{\nu}$，
解得 $\lambda=\frac{hc}{E_{3}-E_{1}}\approx1.03\times10^{-7}\Um$。

评分说明：全题 10 分。（1）3 分，（2）3 分，（3）4 分。

（1）中，列出式（a）的，给 1 分；进一步得出电子动能表达式（b）的，再给 1 分；算出动能值（c）的，再给 1 分。答案（c）中，对能量单位不作限制，但数值和单位必须一致方能得分。

（2）中，正确画出能级图并且标明三个量子数 $n=1$、2、3 和三条光谱线的，给 3 分。每缺一条谱线扣 1 分，缺量子数或量子数不全的扣 1 分，总共最多扣 3 分。

（3）中，算出能级差（d）的，给 1 分；写出频率表达式（e）和波长表达式（f）的，各给 1 分；算出波长（g）的，给 1 分。有效数字只要求三位。
}

\item
%% number: 15
%% paperName: 1982年全国高考第 15 题 16 分
%% typeId: 6
%% type: 计算题
%% chapter: 磁场
%% point: 磁力矩
%% method: 无
%% score: 16
%% degree: 450
%% duplicateId: 0
%% body:
如图 \ref{1982全国15a}，闭合的单匝线圈在匀强磁场中以角速度 $\omega$ 绕中心轴 $OO'$ 逆时针匀速转动。已知：线圈的边长 $ab=cd=l_{1}=0.20\Um$，$bc=da=l_{2}=0.10\Um$，线圈的电阻值 $R=0.050\UO$，角速度 $\omega=300\Urads$；匀强磁场磁感应强度的大小 $B=0.50\UT$，方向与转轴 $OO'$ 垂直。规定当线圈平面与 $B$ 垂直，并且 $ab$ 边在纸面（即过 $OO'$ 轴平行于 $B$ 的平面）前时开始计算线圈的转角 $\theta$。
\twopicture[numA=1, numB=2, labelA=1982全国15a, labelB=1982全国15b, wA=5cm, wB=4cm, align=right]{figs/15a.png}{figs/15b.png}
\begin{enumerate}
	\item
	当 $\theta=\omega t=30^{\circ}$ 时，线圈中感生电动势的大小、方向如何？线圈所受电磁力矩 $M_{\text{磁}}$ 的大小、方向如何？
	\item
	这时，作用在线圈上电磁力的即时功率等于多少？
	\item
	要维持线圈作匀角速转动，除电磁力矩 $M_{\text{磁}}$ 外，还必须另有外力矩 $M_{\text{外}}$ 作用在线圈上。写出 $M_{\text{外}}$ 随时间 $t$ 变化的关系式，并以 $t$ 为横坐标、$M_{\text{外}}$ 为纵坐标画出 $M_{\text{外}}$ 随 $t$ 变化的图线。
\end{enumerate}
%% answer: 
\jdanswer{
\begin{enumerate}
	\item
	$E=1.5\UV$，方向沿 abcd；电磁力矩方向是顺时针的，$M_{\text{磁}}=0.15\UN\cdot\Um$。

	\item
	电磁力的即时功率大小为 $45\UW$。

	\item
	$M_{\text{外}}=0.60\sin^{2}300t$，图线如图所示。
\end{enumerate}
}
%% memo:
\memoanswer{
（1）ab 段和 cd 段切割磁感线的速率为
$v=\omega\frac{l_{2}}{2}=300\times\frac{0.10}{2}\Ums=15\Ums$  ①

线圈内的感生电动势 $e$ 等于 ab 段和 cd 段切割磁感线产生的感生电动势之和，即
$E=2Bl_{1}v\sin\theta=Bl_{1}l_{2}\omega\sin\theta$，
当 $\theta=30^{\circ}$ 时，
$E=0.50\times0.20\times0.10\times300\times0.5\UV=1.5\UV$  ②
方向沿 abcd，

线圈内的电流大小为
$I=\frac{E}{R}=\frac{1.5}{0.050}\UA=30\UA$  ③

\onepicture[align=right, w=4.5cm]{figs/15_an1.png}

ab 段所受安培力 $F_{ab}$ 的方向如图所示，大小为
$F_{ab}=BIl_{1}=30\times0.20\times0.50\UN=3\UN$  ④

安培力 $F_{ab}$ 对于 $OO'$ 轴的力矩方向，在俯视图中是顺时针的，大小为
$M_{ab}=F_{ab}\frac{l_{2}}{2}\sin\theta$。
cd 段所受电磁力矩 $M_{cd}$ 的方向和大小与 ab 段相同。
因此总磁力矩方向是顺时针的，大小为
$M_{\text{磁}}=M_{ab}+M_{cd}=F_{ab}l_{2}\sin\theta=0.15\UN\cdot\Um$  ⑤

（2）ab 段上电磁力的方向与速度方向间的夹角是 $90^{\circ}+\theta$，垂直于速度方向的分力不做功；平行于速度方向的分力为 $F_{ab}\cos\theta$，方向总是与速度相反，所以安培力 $F_{ab}$ 的即时功率大小为
$P_{ab}=|-(F_{ab}\sin\theta)v|$，
故作用在线圈上的电磁力的即时功率为
$P=2P_{ab}=|-2F_{ab}v\sin\theta|=45\UW$。  ⑥

（3）当线圈绕 $OO'$ 轴匀速转动时，线圈所受力矩的代数和等于零，因此，外力矩 $M_{\text{外}}$ 的大小与电磁力矩 $M_{\text{磁}}$ 相等，方向相反，是逆时针的，所以利用 ① 中已经得到的各公式及 $\theta=\omega t$ 可得
$M_{\text{外}}=F_{ab}l_{2}\sin\theta=\frac{B^{2}l_{1}^{2}l_{2}^{2}\omega}{R}\sin^{2}\omega t$  ⑦
代入数据得 $M_{\text{外}}=0.60\sin^{2}300t$  ⑧
以 $t$ 为横坐标和 $M_{\text{外}}$ 为纵坐标的 $M_{\text{外}}-t$ 图线如图所示。

\onepicture[align=right, w=8cm]{\includesvg{figs/15_an2}}

评分说明：全题 16 分。（1）9 分，（2）3 分，（3）4 分。

（1）中，算出速率（a）的，给 1 分；算出电动势（b）的，再给 2 分，得出电动势方向的，再给 1 分。算出电流（c）的，给 1 分；算出 ab 段或 cd 段上电磁力大小（d）的，再给 1 分；得出电磁力方向的，再给 1 分；算出总电磁力矩大小（e）的，再给 1 分；得出总电磁力矩方向的，再给 1 分。在求 $e$ 和 $M_{\text{磁}}$ 时，未分步计算，直接写出 $e=2Bl_{1}v\sin\theta$ 和 $M_{\text{磁}}=Il_{1}l_{2}B\sin\theta$ 的，不扣分。

（2）中，算出电磁力即时功率（f）的，给 3 分。符号错误的，扣 1 分。按其他方法计算，例如按 $P=IE$ 或 $P=I^{2}R$ 计算的，同样给分。

（3）中，得出 $M_{\text{外}}$ 表达式（g）或（h）的，给 2 分；作出 $M_{\text{外}}-t$ 曲线的，再给 2 分。在图上，将 $2\pi/\omega$ 标作 $\pi/150$ 以及将 $(Bl_{1}l_{2})^{2}\omega/R$ 标作 0.6 的，同样给分。曲线上只要 $t=0$、$\pi/2\omega$、$\pi/\omega$、$3\pi/2\omega$、$2\pi/\omega$ 各时刻的 $M_{\text{外}}$ 值正确，形状虽与正弦平方曲线稍有出入，不扣分。所得 $M_{\text{外}}$ 表达式与（g）式或（h）式差一负号，并且曲线全在横轴下方的，不扣分。
}

\item
%% number: 16
%% paperName: 1982年全国高考第 16 题 12 分
%% typeId: 6
%% type: 计算题
%% chapter: 运动和力
%% point: 动量和能量
%% method: 无
%% score: 12
%% degree: 350
%% duplicateId: 0
%% body:
在一原子反应堆中，用石墨（碳）作减速剂使快中子减速。已知碳核的质量是中子的 12 倍。假设把中子与碳核的每次碰撞都看作是弹性正碰，而且认为碰撞前碳核都是静止的。
\begin{enumerate}
	\item
	设碰撞前中子的动能是 $E_{0}$，问经过一次碰撞中子损失的能量是多少？
	\item
	至少经过多少次碰撞，中子的动能才能小于 $10^{-6}E_{0}$？（$\lg13=1.114$，$\lg11=1.041$。）
\end{enumerate}
%% answer: 
\jdanswer{
\begin{enumerate}
	\item
	$\Delta E=\frac{48}{169}E_{0}$

	\item
	42 次
\end{enumerate}
}
%% memo:
\memoanswer{
（1）设中子和碳核的质量分别为 $m$ 和 $M$，碰撞前中子的速度为 $v_{0}$，碰撞后中子和碳核的速度分别为 $v$ 和 $V$，由动量守恒得
$mv_{0}=mv+MV$  ①，
由能量守恒得
$\frac{1}{2}mv_{0}^{2}=\frac{1}{2}mv^{2}+\frac{1}{2}MV^{2}$  ②，
由 ①② 两式可解出 $v_{0}+v=V$  ③。
将上式代入 ① 中，得
$v=\frac{m-M}{m+M}v_{0}$。
已知 $M=12m$，代入 ③ 中可得
$v=-\frac{11}{13}v_{0}$  ④。
因此在碰撞过程中中子损失的能量为
$\Delta E=\frac{1}{2}mv_{0}^{2}-\frac{1}{2}mv^{2}$  ⑤
$=E_{0}\left(1-\frac{121}{169}\right)=\frac{48}{169}E_{0}$  ⑥。

（2）设 $E_{1},E_{2},\ldots,E_{n}$ 分别表示中子在第 1 次、第 2 次……第 $n$ 次碰撞后的动能。由 ④ 可得
$E_{1}=\left(\frac{11}{13}\right)^{2}E_{0}$  ⑦。
同理可得
$E_{2}=\left(\frac{11}{13}\right)^{2}E_{1}=\left(\frac{11}{13}\right)^{4}E_{0}$，
$\cdots$
$E_{n}=\left(\frac{11}{13}\right)^{2n}E_{0}$  ⑧，
已知 $E_{n}=10^{-6}E_{0}$，代入上式可得
$10^{-6}E_{0}=\left(\frac{11}{13}\right)^{2n}E_{0}$，即 $\left(\frac{11}{13}\right)^{2n}=10^{-6}$  ⑨，
取对数可得
$2n(\lg13-\lg11)=6$。
将 $\lg13=1.114$，$\lg11=1.041$ 代入上式，即得
$n=\frac{3}{0.073}=41.1$，
因此，至少需经过 42 次碰撞，中子的动能才能小于 $10^{-6}E_{0}$。

评分说明：全题 12 分。（1）6 分，（2）6 分。

（1）中，列出（a）和（b）的，各给 1 分；从（a）、（b）解出（c）的，再给 2 分。不写出（a）和（b）直接写出（c）的，只给 1 分。得出（e）的，再给 1 分；求出结果（f）的，再给 1 分。

（2）中，列出（h）的，给 4 分；只写出（g）的，给 1 分。将 $E_{n}=10^{-6}E_{0}$ 代入（h），并得出（i）的，再给 1 分。由（i）算出最后结果的，再给 1 分。
}

\end{enumerate}

\end{document}
